\section{Discrete calculation of the measure on the inverse limit: differences, carry cocycle, ultrametric, reconstruction of \texorpdfstring{\(\mathbb R\)}{R}, and genealogical disintegration}
\label{sec:complejo-discreto-medida-rev2}
\index{measure!discrete complex}
\index{path!integration}
\index{carry!cocycle}

The distinction between counting and measuring admits an operator formulation. Counting assigns a cardinality to a collection of marks. Measuring evaluates relations among marks by means of a local rule, integrates them along a path, and preserves the chart in which the result is published. The minimal unit of that process is a difference; its composition produces length, circulation, and, when a modular reduction intervenes, carry. This architecture places the measure prior to its real coordinate: the published scalar is the contraction of a calculation possessing a domain, orientation, and memory.

\subsection{Marks and difference operator}

Let \(P_n\) be the path with ordered vertices
\(V(P_n)=\{v_0,\ldots,v_{n-1}\}\) and oriented edges
\(e_j=(v_{j-1},v_j)\), \(1\leq j\leq n-1\). For a commutative ring
\(A\), we define
\[
 C^0(P_n;A)=A^{V(P_n)},
 \qquad
 C^1(P_n;A)=A^{E(P_n)},
\]
and the difference operator
\begin{equation}
 d:C^0(P_n;A)\longrightarrow C^1(P_n;A),
 \qquad
 (df)(e_j)=f(v_j)-f(v_{j-1}).
 \label{eq:diferencia-cero-uno-rev2}
\end{equation}

\begin{theorem}[Range of the difference operator on a path]
\label{thm:rango-diferencias-camino-rev2}
If \(A\) is a field, then
\[
 \ker d=A\mathbf 1,
 \qquad
 \operatorname{rank}d=n-1.
\]
Consequently, \(n\) connected marks determine \(n-1\)
independent differences and a unique constant mode.
\end{theorem}

\begin{proof}
The equality \(df=0\) successively implies
\(f(v_0)=f(v_1)=\cdots=f(v_{n-1})\), so that the kernel is formed by
the constant functions and has dimension one. The rank--nullity theorem
gives \(\operatorname{rank}d=n-1\). Equivalently, every
assignment \((a_1,\ldots,a_{n-1})\) to the edges is integrated uniquely
once \(f(v_0)\) has been fixed:
\[
 f(v_k)=f(v_0)+\sum_{j=1}^{k}a_j.
\]
\end{proof}

The constant mode expresses freedom of choosing the origin. Differences
contain relational information; the initial mark supplies the section
reconstructing positions. In the inaugural HMT cell, nine ordered marks
produce eight interior differences. Periodic closure adds the ninth
edge and transforms the path into a cycle; the two-dimensional completion
of those intervals will be realized exactly in the APP construction.

\subsection{Positive length and oriented integral}

Two evaluations of a route must remain separate. A discrete metric is a positive cochain.
\[
 \ell:E(P_n)\longrightarrow\mathbb R_{>0}.
\]
For an unoriented path \(|\gamma|\), its length is
\begin{equation}
 L_\ell(\gamma)=\sum_{e\in|\gamma|}m_\gamma(e)\,\ell(e),
 \label{eq:longitud-positiva-camino-rev2}
\end{equation}
where \(m_\gamma(e)\) counts the traversals of the edge. An oriented
cochain \(\omega\in C^1(P_n;A)\), by contrast, is integrated with sign:
\begin{equation}
 \int_\gamma\omega
 :=\sum_{k=1}^{r}\varepsilon_k\omega(e_{i_k}),
 \qquad \varepsilon_k\in\{-1,+1\}.
 \label{eq:integral-cocadena-camino-rev2}
\end{equation}
Length sums geometric cost; the oriented integral sums transport. The
first is invariant under route inversion, and the second changes sign.

\begin{proposition}[Discrete Fundamental Theorem]
\label{prop:teorema-fundamental-discreto-rev2}
For \(f\in C^0(P_n;A)\) and every path orientada \(\gamma:a\to b\),
\[
 \int_\gamma df=f(b)-f(a).
\]
In particular, the integral of an exact cocycle over a cycle is zero.
\end{proposition}

\begin{proof}
The sum of the consecutive differences telescopes. Each interior vertex
appears once with positive sign and once with negative sign; only
the endpoints remain.
\end{proof}

The proposition distinguishes two levels of closure. The group
\(H_1(C_n;\mathbb Z)\cong\mathbb Z\) provides the topological carrier of the
turn; a holonomy further requires a transport. If
\(T:\operatorname{Path}(C_n)\to\mathcal C\) is a functor toward a category
of automorphisms, the holonomy of a cycle \(\gamma\) is
\begin{equation}
 \operatorname{Hol}_T(\gamma)=T(e_r)\cdots T(e_1).
 \label{eq:holonomia-transporte-ciclo-rev2}
\end{equation}
The cycle contributes the traversal class; \(T\) contributes the discrete connection that
turns that class into an endomorphism. This distinction will be essential when
the visible-phase return recovers the phase and the transport conserves a different
boundary state.

\subsection{Carry as a composition cocycle}

Modular reduction publishes a class, and a section chooses representatives.
Let \(Q=\mathbb Z/9\mathbb Z\), and let
\(s_0:Q\to\{0,\ldots,8\}\subset\mathbb Z\) be the ordinary residue section.
We define
\begin{equation}
 \kappa_9(a,b)
 :=\frac{s_0(a)+s_0(b)-s_0(a+b)}{9}\in\mathbb Z.
 \label{eq:cociclo-acarreo-nonadico-rev2}
\end{equation}
Then
\[
 s_0(a)+s_0(b)=s_0(a+b)+9\kappa_9(a,b).
\]

\begin{proposition}[Cocycle Identity of the Carry]
\label{prop:acarreo-dos-cociclo-rev2}
The map \(\kappa_9:Q\times Q\to\mathbb Z\) satisfies
\begin{equation}
 \kappa_9(a,b)+\kappa_9(a+b,c)
 =\kappa_9(b,c)+\kappa_9(a,b+c).
 \label{eq:identidad-dos-cociclo-rev2}
\end{equation}
It is, therefore, a normalized \(2\)-cocycle associated with the chosen section.
\end{proposition}

\begin{proof}
Both sides are expanded using
\eqref{eq:cociclo-acarreo-nonadico-rev2}. After multiplication by nine, both
are equal to
\[
 s_0(a)+s_0(b)+s_0(c)-s_0(a+b+c).
\]
Normalization is obtained from \(s_0(0)=0\).
\end{proof}

The cocyclic identity expresses the associativity of lifting: grouping first
\((a+b)+c\) or \(a+(b+c)\) produces the same visible class and the same total
carry. In a category of routes, the same structure takes the form
\[
 \widetilde g\,\widetilde h
 =\kappa(g,h)\,\widetilde{gh},
\]
where \(\widetilde g,\widetilde h\) are enriched transports lifting
visible morphisms. The carry is thus a perfectly controlled composition
defect before receiving any geometric interpretation as
curvature or torsion.

\subsection{Compatibility of the nonadic, ternary, and decimal-completion charts in base \texorpdfstring{\(10^3\)}{10^3}}
\label{subsec:compatibilidad-lectores-9-3-1000-rev3}

The residual readers that intervene in TPK are not distinct names of
a single operation. Let
\[
 R_9:=\mathbb Z/9\mathbb Z,
 \qquad
 r_9:\mathbb Z\longrightarrow R_9,
 \qquad
 r_3:\mathbb Z\longrightarrow\mathbb F_3
\]
the canonical reductions, and let
\[
 \pi_3:R_9\longrightarrow\mathbb F_3,
 \qquad \pi_3(\overline n)=n\bmod3.
\]
The square of reductions commutes:
\begin{equation}
\begin{CD}
 \mathbb Z @>{r_9}>> R_9\\
 @V{r_3}VV @VV{\pi_3}V\\
 \mathbb F_3 @= \mathbb F_3 .
\end{CD}
\qquad
 r_3=\pi_3\circ r_9.
\label{eq:diagrama-reducciones-9-3-rev3}
\end{equation}
The positive section
\[
 s_+:R_9\longrightarrow\{1,\ldots,9\},
 \qquad
 s_+(\overline0)=9,
 \qquad
 s_+(\overline k)=k\quad(1\le k\le8),
\]
produce, for \(n>0\), the root digital positive
\[
 \operatorname{dr}_9^+(n)=s_+(r_9(n)).
\]
This section is not a homomorphism: it publishes the null class by means of \(9\) and
forgets the quotient. The residual section \(s_0(\overline0)=0\) publishes that
same class by means of \(0\). Therefore, the visible change \(9\mapsto0\) is a
change of section, and the equality
\[
 9=s_0(\overline0)+9\cdot1
\]
preserves the carry \(1\) as quotient. The topological successor of the
positive cycle is, by contrast, \(9\mapsto1\), with
\[
 10=s_+(\overline1)+9\cdot1.
\]
Thus the closure \(P_9\to C_9\), the change of chart and the increment
of memory are separated; none of the three is replaced by the other two.

The coordinate of the radical belongs to another diagram. For
\(\mathfrak m=(3)\subset R_9\), the application
\[
 \iota_{\mathfrak m}:\mathbb F_3\longrightarrow\mathfrak m,
 \qquad \overline a\longmapsto3a,
\]
is an isomorphism of vector spaces over \(\mathbb F_3\), with inverse
\begin{equation}
 \eta_{\mathfrak m}:\mathfrak m\longrightarrow\mathbb F_3,
 \qquad
 \eta_{\mathfrak m}(3a)=a\bmod3.
 \label{eq:coordenada-radical-compatibilidad-rev3}
\end{equation}
It is not a ring isomorphism. Nor does it coincide with direct reduction:
for \(3a\in\mathfrak m\), one has \(\pi_3(3a)=0\), whereas \(\eta_{\mathfrak m}(3a)=a\bmod3\). The first
reading publishes the residue class; the second publishes the internal
coordinate of the radical.

Finally, let \(u_j\in\{0,\ldots,999\}\) and
\[
 D_N=1000D_{N-1}+u_N,
 \qquad D_0=0,
\]
the concatenation of blocks in base \(10^3\). Since \(10^3\equiv1\pmod9\),
\begin{equation}
 r_9(D_N)=r_9\!\left(\sum_{j=1}^{N}u_j\right),
 \qquad
 r_3(D_N)=r_3\!\left(\sum_{j=1}^{N}u_j\right).
 \label{eq:compatibilidad-base-mil-reducciones-rev3}
\end{equation}
The basis \(10^3\) preserves the order of the blocks; the reductions preserve only their accumulated class. The digital root adds a positive section; the radical coordinate recovers an internal direction of \((3)\). This hierarchy fixes what information each reader publishes and what quotient, order, or chart TPK memory must retain in order to reconstruct the antecedent.

\subsection{Graded system of measurement and non-contracted recording}

The three previous levels are organized in a graded system:
\[
 \boxed{
 \begin{array}{ccl}
 \text{degree }0&:&\text{values or states on marks},\\
 \text{degree }1&:&\text{differences and transports along edges},\\
 \text{degree }2&:&\text{composition and carry cocycles}.
 \end{array}}
\]
Degrees zero and one form the cellular complex
\(C^0\xrightarrow d C^1\); degree two lies in the nerve complex of the
path category. Their union constitutes the \emph{discrete
measure complex}: an organizing structure that preserves the distinction between
cellular geometry and categorical memory.

For an enriched route
\(\Gamma=(x_0\xrightarrow{e_1}x_1\to\cdots\xrightarrow{e_r}x_r)\), the
register is the uncontracted integral
\begin{equation}
 \Lambda(\Gamma)
 =\bigl((x_{k-1},e_k,x_k),\omega_k,\kappa_k,\sigma_k,b_k,m_k\bigr)_{k=1}^{r},
 \label{eq:ledger-integral-no-contraida-rev2}
\end{equation}
where \(\omega_k\) is the local reading, \(\kappa_k\) the carry, \(\sigma_k\)
the orientation or partial signature, \(b_k\) the boundary datum, and \(m_k\) the
preserved memory. An observable is a typed map
\[
 \mathcal O=F_{\mathcal O}\circ\Lambda.
\]
Evaluation may publish a sum, a residual class, a signature, or a scalar without converting those outputs into complete states. Concatenation of routes concatenates registers and adds the cocycle at the common boundary; the uncontracted register therefore preserves exactly the information required to recompose the history that an isolated coordinate contracts.

\begin{figure}[tbp]
\centering
\begin{tikzpicture}[
  >=Latex,
  every node/.style={font=\small,align=center},
  object/.style={draw,rounded corners=2pt,inner sep=5pt,text width=3.00cm,
                 minimum height=1.25cm},
  arrow/.style={->,semithick},
  maplabel/.style={font=\footnotesize,fill=white,inner sep=1.5pt}
]
\node[object] (marks) at (0,0)
  {marks and states\\\(C^0\)};
\node[object] (edges) at (5.25,0)
  {differences and transport\\\(C^1\)};
\node[object] (paths) at (10.50,0)
  {paths and cycles\\integration and holonomy};
\node[object] (memory) at (0,-2.35)
  {composition cocycle\\carry and boundary};
\node[object] (section) at (5.25,-2.35)
  {uncontracted record\\compatible section};
\node[object] (outputs) at (10.50,-2.35)
  {evaluations typed\\value, signature or symmetry};
\draw[arrow] (marks) -- node[maplabel,above]{\(d\)} (edges);
\draw[arrow] (edges) -- node[maplabel,above]{\(\int_\gamma\)} (paths);
\draw[arrow] (paths) -- node[maplabel,right]{transporte TPK} (outputs);
\draw[arrow] (memory) -- node[maplabel,above]{composition} (section);
\draw[arrow] (section) -- node[maplabel,above]{contraction} (outputs);
\draw[arrow] (edges) -- node[maplabel,left]{lifting defect} (memory);
\draw[arrow] (paths) -- node[maplabel,right]{history} (section);
\end{tikzpicture}
\caption{Architecture of the discrete measure complex. The difference
operator, path integral, and composition cocycle occupy
distinct domains; TPK assembles them into a compatible section before
a reader publishes a value or a symmetry.}
\label{fig:complejo-discreto-medida-rev3}
\end{figure}

\subsection{Reconstruction of the continuum from the inverse system of compatible cylinders}

Prolongation to arbitrary depth is formulated by means of an inverse system.
Let \(S_n\) be nonempty finite sets of enriched prefixes and
\(r_{n+1,n}:S_{n+1}\to S_n\) surjective restriction maps. Their
space of compatible sections is
\[
 \Omega_\infty=\varprojlim_n S_n
 =\{(x_n)_n:r_{n+1,n}(x_{n+1})=x_n\}.
\]
Each \(S_n\) is endowed with the discrete topology and \(\Omega_\infty\) with the
subspace topology of the product.

\subsection{Cylinders, ultrametric and Archimedean return with memory conservation}

For \(a\in S_n\), the prefix cylinder \(a\) is
\begin{equation}
 [a]_n:=\{x\in\Omega_\infty:x_n=a\}.
 \label{eq:cilindro-genealogico-rev3}
\end{equation}
Two cylinders at the same level are disjoint, and every cylinder at level
\(n+1\) is contained in the cylinder of its restriction. This clopen
basis publishes the depth shared by two sections without
contracting their memory to a real coordinate.

For \(x\neq y\), let
\[
 N(x,y):=\min\{n\geq0:x_n\neq y_n\}.
\]
Fijado \(0<\vartheta<1\), definimos
\begin{equation}
 d_\vartheta(x,y):=\vartheta^{N(x,y)},
 \qquad d_\vartheta(x,x):=0.
 \label{eq:ultrametrica-genealogica-rev3}
\end{equation}

\begin{proposition}[First Divergence Ultrametric]
\label{prop:ultrametrica-genealogica-rev3}
The function \(d_\vartheta\) is an ultrametric, and its balls are precisely
the cylinders, with the index adjustment corresponding to the initial level:
\[
 d_\vartheta(x,z)
 \leq
 \max\{d_\vartheta(x,y),d_\vartheta(y,z)\}.
\]
Therefore, two sections that publish the same value can remain at
positive distance when their sheet, carry, boundary or
history of prolongation diverges.
\end{proposition}

\begin{proof}
If \(x\) and \(y\) agree through level \(n\), and \(y\) and \(z\) do as well,
then \(x\) and \(z\) agree through that level. The first divergence of the
first and third states therefore occurs no earlier than the minimum of
the other two. Since \(0<\vartheta<1\), raising \(\vartheta\) to those
exponents yields the strong inequality. The description of the balls follows
directly from fixing a common prefix.
\end{proof}

\begin{theorem}[Compactness and completeness of the genealogical space]
\label{thm:compacidad-ultrametrica-genealogica-rev3}
If each \(S_n\) is finite and the bonding maps are surjective, then
\(\Omega_\infty\), with \(d_\vartheta\), is nonempty, compact, and complete.
If, in addition,
\begin{equation}
 \forall x\in\Omega_\infty\ \forall n\ \exists y\neq x:
 y_n=x_n,
 \label{eq:perfeccion-genealogica-rev3}
\end{equation}
has no isolated points. In that regime it is a compact metrizable,
zero-dimensional, and perfect space; by the classical characterization, it is homeomorphic
to Cantor space.
\end{theorem}

\begin{proof}
Surjectivity and finite branching provide a section by König's lemma. The product of the finite discrete sets is compact, and the equations \(r_{n+1,n}(x_{n+1})=x_n\) define a closed subspace. The induced topology coincides with the cylinder topology and, by \cref{prop:ultrametrica-genealogica-rev3}, with that of \(d_\vartheta\). Every compact metric space is complete. The condition \eqref{eq:perfeccion-genealogica-rev3} guarantees that every cylinder contains a section distinct from its center, so no singleton is open.
\end{proof}

Visible periodicity and state return belong to distinct
maps. Let \(p:E\twoheadrightarrow B\) be a projection, \(T:E\to E\) an
enriched transport, and \(f:B\to B\) the published evolution, with
\(p\circ T=f\circ p\). If \(f^9=I_B\), then
\begin{equation}
 p\circ T^9=p,
 \label{eq:retorno-visible-fibra-rev3}
\end{equation}
so that \(T^9\) preserves each fiber of \(p\), but need not be
the identity on it. In the nonadic coordinate
\[
 (q,r)\longmapsto
 \begin{cases}
  (q,r+1),&r<9,\\
  (q+1,1),&r=9,
 \end{cases}
\]
one revolution recovers the position and increases the height. Under CT108
periodicity, nine iterations advance one dodecaphase phase, and one hundred eight
recover position and phase; the enriched state may still preserve
different sheet, boundary, orientation, and genealogy. This is the formal
distinction between phase return and state return.

\begin{theorem}[Reconstruction by Compatible Sections]
\label{thm:reconstruccion-continuo-secciones-rev2}
The space \(\Omega_\infty\) is nonempty and compact. Suppose that to each
prefix \(x_n\in S_n\) there is assigned a compact interval
\(I_n(x_n)\subset\mathbb R\) such that, for every section
\(x=(x_n)_n\),
\[
 I_{n+1}(x_{n+1})\subseteq I_n(x_n),
 \qquad
 \operatorname{diam}I_n(x_n)\longrightarrow0.
\]
Then there exists a unique value
\begin{equation}
 \operatorname{ev}_{\mathbb R}(x)
 \in\bigcap_{n\geq0}I_n(x_n).
 \label{eq:evaluacion-seccion-continua-rev2}
\end{equation}
If the interval assignments are locally constant at each level, the
evaluation \(\operatorname{ev}_{\mathbb R}:\Omega_\infty\to\mathbb R\) is
continuous.
\end{theorem}

\begin{proof}
Surjectivity permits compatible prefixes to be constructed inductively; compactness of the finite sets and the finite-intersection property yield a global section. The product \(\prod_nS_n\) is compact, and \(\Omega_\infty\) is closed, since every compatibility equation defines a closed set. For a fixed section, the intervals are compact, nested, and of diameter tending to zero; the nested-interval theorem yields a unique common point. Given \(\varepsilon>0\), one chooses a level whose interval has diameter less than \(\varepsilon\). Sections sharing that prefix have evaluations within the same interval, which proves continuity.
\end{proof}

The real line thus appears as the codomain of a contraction. Two sections
may share a value while preserving different records; the relation
\[
 x\sim_{\mathbb R}y
 \quad\Longleftrightarrow\quad
 \operatorname{ev}_{\mathbb R}(x)=\operatorname{ev}_{\mathbb R}(y)
\]
forms the Archimedean quotient. The genealogy remains in \(\Omega_\infty\), whereas the coordinate resides in its real image. The continuum problem therefore receives a constructive answer within the domain covered by the declared prefix systems: continuity of evaluation and memory of the section are coordinated properties, not competing descriptions.

\subsection{Coherent coverage and reconstruction of the Archimedean quotient}

The preceding evaluation assigns a value to each section. To prove that the
visible space is completely reconstructed, it is also necessary that
refinement lose no points. Let \(K\subset\mathbb R\) be compact, and suppose
that the intervals of the preceding theorem satisfy
\begin{align}
 K&=\bigcup_{a\in S_0} I_0(a),
 \label{eq:cobertura-inicial-continuo-rev3}\\
 I_n(a)&=\bigcup_{\substack{b\in S_{n+1}\\r_{n+1,n}(b)=a}}I_{n+1}(b)
 \qquad(a\in S_n).
 \label{eq:cobertura-hijos-continuo-rev3}
\end{align}
These identities define a \emph{coherent covering}: every visible
compact decomposes exactly into the compacts of its prolongations.

\begin{theorem}[Genealogical Reconstruction of a Compact Set]
\label{thm:cociente-genealogico-compacto-rev3}
Suppose that:
\begin{enumerate}
\item the sets \(S_n\) are finite and nonempty, and the bonding maps
      \(r_{n+1,n}\) are surjective;
\item the intervals are compact, nonempty, nested, and uniformly
      contractive;
\item
      \eqref{eq:cobertura-inicial-continuo-rev3} and
      \eqref{eq:cobertura-hijos-continuo-rev3} hold.
\end{enumerate}
Then \(\operatorname{ev}_{\mathbb R}(\Omega_\infty)=K\), and the map
induced
\begin{equation}
 \overline{\operatorname{ev}}_{\mathbb R}:
 \Omega_\infty/\!\sim_{\mathbb R}\xrightarrow{\ \cong\ }K
 \label{eq:homeomorfismo-cociente-continuo-rev3}
\end{equation}
is a homeomorphism. The fiber over \(y\in K\) is the complete set of
compatible sections that publish that same value.
\end{theorem}

\begin{proof}
Compactness and continuity follow from
\cref{thm:reconstruccion-continuo-secciones-rev2}. With \(y\in K\) fixed, let
\[
 T_n(y):=\{a\in S_n:y\in I_n(a)\}.
\]
Coherent coverage makes every \(T_n(y)\) nonempty and connects each
state to a prolongation. The resulting tree is finitely branching and
has vertices at every depth; by the infinite-branch lemma it admits a
compatible chain \((x_n)_n\). The point \(y\) belongs to every
\(I_n(x_n)\), and contraction forces
\(\operatorname{ev}_{\mathbb R}(x)=y\). Therefore, the evaluation is
surjective. The quotient map is continuous and bijective; since its
domain is compact and \(K\) is Hausdorff, it is a homeomorphism.
\end{proof}

The theorem distinguishes two assertions. The contraction produces a unique value
for each genealogy; the coherent covering proves that all values
of the compact set appear. A check at finite depth certifies the
examined levels, but it does not replace any of these infinite hypotheses.

\begin{proposition}[Genealogical Atlas of the Real Line]
\label{prop:atlas-genealogico-recta-real-rev3}
For each \(m\geq1\), let \(K_m=[-m,m]\), and suppose that a system is given that
satisfies \cref{thm:cociente-genealogico-compacto-rev3}, together with
closed inclusions
\[
 \jmath_m:\Omega_\infty^{(m)}\hookrightarrow\Omega_\infty^{(m+1)}
\]
that intertwine the evaluations and the inclusion
\(K_m\hookrightarrow K_{m+1}\). Then the Archimedean quotient of the
strict colimit of the sections is
\[
 \varinjlim_m K_m=\bigcup_{m\geq1}[-m,m]=\mathbb R,
\]
with its usual topology, while the fibers of the colimit preserve the genealogies of each value.
\end{proposition}

\begin{proof}
Each piece has quotient \(K_m\) by
\eqref{eq:homeomorfismo-cociente-continuo-rev3}. The declared inclusions
make these quotients compatible. The final topology of the exhaustion by
compact intervals coincides with the usual topology of \(\mathbb R\): a
set is open exactly when its intersection with each \(K_m\) is
open in \(K_m\).
\end{proof}

\subsection{Algebraic descent of operations}

Reconstructing the visible space does not suffice to reconstruct its algebra. Let
\(\star\) be a continuous partial operation on \(\Omega_\infty\), with domain
\(D_\star\) saturated by the fibers of
\(\operatorname{ev}_{\mathbb R}\times\operatorname{ev}_{\mathbb R}\).

\begin{theorem}[Fiberwise Descent Criterion]
\label{thm:descenso-operaciones-continuo-rev3}
There exists a visible operation \(\overline\star\) that satisfies
\[
 \operatorname{ev}_{\mathbb R}(x\star y)
 =\operatorname{ev}_{\mathbb R}(x)\,\overline\star\,
  \operatorname{ev}_{\mathbb R}(y)
\]
if and only if the value of the left-hand side is constant upon replacing
\(x\) or \(y\) by any other representative of its respective fiber.
If two descended continuous operations coincide with addition and multiplication
on dense subsets of their visible domains, they coincide with them throughout
the domain.
\end{theorem}

\begin{proof}
Constancy on fibers is precisely the universal property of the quotient: it
permits \(\overline\star\) to be defined independently of the representative and
is necessary for that definition to be meaningful. The final assertion
follows from continuity and density.
\end{proof}

APP conserves the datum required to address this descent: the sum and
product leaves act on the same support, but maintain different quotients, carries, and
fibers. The visible operation can descend even though the memory of
its representatives remains distinct.

\subsection{Projective measure and conservation of the unit}

Cylinder topology organizes proximity and truncation; a path theory also
requires weights. Let \(\mu_n\) be a probability measure on
\(S_n\). The family is projectively consistent when
\begin{equation}
 \mu_n(a)=
 \sum_{\substack{b\in S_{n+1}\\r_{n+1,n}(b)=a}}\mu_{n+1}(b).
 \label{eq:consistencia-medida-cilindrica-rev3}
\end{equation}
The equality expresses the local conservation of the unit: the weight of a
truncated history is the sum of the weights of all its prolongations.

\begin{theorem}[Genealogical Measure on the Inverse Limit]
\label{thm:medida-genealogica-limite-rev3}
Every family of finite probabilities satisfying
\eqref{eq:consistencia-medida-cilindrica-rev3} determines a unique
Borel probability \(\mu_\infty\) on \(\Omega_\infty\) such that
\begin{equation}
 \mu_\infty([a]_n)=\mu_n(a)
 \label{eq:medida-cilindro-rev3}
\end{equation}
for every prefix cylinder \(a\in S_n\).
\end{theorem}

\begin{proof}
Finite unions of cylinders form an algebra. Any one of them can be refined into a disjoint union at a common level, and \eqref{eq:consistencia-medida-cilindrica-rev3} proves that the sum of their weights does not depend on the chosen level. A premeasure is thereby obtained. Since the cylinders are both open and closed and \(\Omega_\infty\) is compact, a decreasing sequence of elements of the algebra with empty intersection already reaches the empty set at some level. The premeasure is therefore continuous at the empty set and extends uniquely to the \(\sigma\)-algebra of Borel sets.
\end{proof}

The evaluation transports that measure to the visible compact,
\begin{equation}
 \nu=(\operatorname{ev}_{\mathbb R})_*\mu_\infty,
 \qquad
 \nu(B)=\mu_\infty\bigl(\operatorname{ev}_{\mathbb R}^{-1}(B)\bigr).
 \label{eq:pushforward-medida-visible-rev3}
\end{equation}
The visible probability adds the genealogies of a single fiber without
erasing them from the total space.

\begin{proposition}[Value Decomposition and Genealogical Multiplicity]
\label{prop:desintegracion-valores-genealogicos-rev3}
If \(K\) is a standard Borel space, then for \(\nu\)-almost every
\(y\in K\) there exists a conditional probability \(\mu_y\) supported on
\[
 \mathfrak G_y:=\operatorname{ev}_{\mathbb R}^{-1}(y)
\]
such that, for every integrable function \(F\),
\begin{equation}
 \int_{\Omega_\infty}F\,\mathrm d\mu_\infty
 =\int_K
   \left(\int_{\mathfrak G_y}F\,\mathrm d\mu_y\right)
   \mathrm d\nu(y).
 \label{eq:desintegracion-fibras-visibles-rev3}
\end{equation}
Thus, the visible coordinate and the distribution of histories that realize it
are distinct typed data: an observable that factors through
\(\operatorname{ev}_{\mathbb R}\) reads only \(\nu\), whereas a memory
reader can distinguish the measures \(\mu_y\) within a single fibre.
\end{proposition}

\begin{proposition}[Compatible cylinder conditioning]
\label{prop:condicionamiento-cilindrico-compatible-rev3}
Let \(A\subseteq\Omega_\infty\) be a union of cylinders of level \(m\), with
\(\mu_\infty(A)>0\). For every \(n\geq m\), the conditioned distribution
on \(S_n\) is obtained by suppressing the states whose truncation does not
belong to \(A\) and renormalizing the remaining weights. The resulting
distributions continue to satisfy projective consistency
\eqref{eq:consistencia-medida-cilindrica-rev3}.
\end{proposition}

\begin{proof}
This is Bayes' rule on the finite sets at level \(n\). The sum of the conditional weights of the prolongations equals the conditional weight of the parent because intersection with \(A\) distributes over the disjoint union of its prolongations. Conditioning eliminates incompatible futures and redistributes the unit among the survivors; it neither creates new histories nor removes from the total space those lying outside the event being read.
\end{proof}

This separation provides the exact basis for distinguishing preparation,
reading, conditioning, and history in the quantum realization.

\begin{statusbox}[title={Mathematical force and architectural function}] The rank \(n-1\), the discrete fundamental theorem, the carry cocycle identity, the reconstruction theorems, and the projective measure are exact results in the declared domains. Their HMT application uses as its starting structure the states, restrictions, intervals, and transports constructed by APP--TRIT--TPK. The discrete measure complex is an organizing formalization: it does not identify the degree-two categorical cocycle with a geometric cell until the corresponding map is provided. This separation allows the common language subsequently to govern holonomy, the sum over paths, and the transition from number to particle without anticipating their physical realizations.
\end{statusbox}
